\chapter{Mathematical Proofs}
\label{app:proofs}

This appendix collects the full proofs of the formal results stated in Chapter~\ref{ch:methodology}. Each section below corresponds to a single proposition, lemma, or corollary, and may be read independently of the others. Throughout the appendix, expectations are taken with respect to the on-policy distribution induced by the current policy $\pi_\theta$ unless stated otherwise; standard regularity conditions (existence of first and second moments, measurability of $\beta_\psi$, and a state-visitation distribution that admits a density) are assumed.

\section{Proof of Proposition 3.1: Target Recovery}
\label{app:proof_target}

\begin{proof}
We work with the population form of the learning-progress loss in Equation~\eqref{eq:progress_loss}:
\begin{equation*}
L_\text{progress}(\psi) \;=\; \mathbb{E}\!\left[ \bigl( \beta_\psi(s_t) - \tau(s_t) \bigr)^2 \right],
\end{equation*}
where the expectation is over the on-policy state-visitation distribution. Setting $\lambda_\text{reg} = 0$ in Equation~\eqref{eq:beta_loss} reduces the Beta-network objective to $L_\beta(\psi) = \gamma_p\, L_\text{progress}(\psi)$, which is minimized at the same $\psi$ as $L_\text{progress}$.

Assume the Beta Network has sufficient capacity to represent any measurable function $f : \mathcal{S} \to [\beta_\text{min}, \beta_\text{max}]$. Then minimizing $L_\text{progress}$ over $\psi$ is equivalent to minimizing
\begin{equation*}
\mathcal{J}(f) \;=\; \mathbb{E}\!\left[ \bigl( f(s_t) - \tau(s_t) \bigr)^2 \right]
\end{equation*}
over the function class. Conditioning on $s_t$ and applying the law of total expectation,
\begin{equation*}
\mathcal{J}(f) \;=\; \mathbb{E}\!\left[ \mathbb{E}\bigl[ (f(s_t) - \tau(s_t))^2 \,\big|\, s_t \bigr] \right].
\end{equation*}
Because $f(s_t)$ is constant given $s_t$, the inner expectation is the conditional mean-squared error of $f(s_t)$ as a predictor of $\tau(s_t)$. By the standard $L^2$-projection argument, this inner expectation is uniquely minimized (over choices of $f(s_t)$) by the conditional mean
\begin{equation*}
f^\star(s) \;=\; \mathbb{E}\!\left[\, \tau(s_t) \,\big|\, s_t = s \right].
\end{equation*}
Since $\tau(s_t) \in [\beta_\text{min}, \beta_\text{max}]$ almost surely (by Equation~\eqref{eq:beta_target} together with $\beta_0 \in [\beta_\text{min}, \beta_\text{max}]$), the conditional mean $f^\star(s)$ also lies in $[\beta_\text{min}, \beta_\text{max}]$, and is therefore representable by the Beta Network. Setting $\beta_{\psi^\star}(s) = f^\star(s)$ yields a global minimizer of $L_\beta$, and the strict convexity of the squared loss makes it unique on the support of the state-visitation distribution.
\end{proof}

\section{Proof of Proposition 3.2: Gradient Direction}
\label{app:proof_gradient}

\begin{proof}
Fix a state $s \in \mathcal{S}$ visited with positive probability, and treat $\beta_\psi(s)$ as a free real-valued parameter for the purpose of computing the partial derivative. The population form of the progress loss can be written as
\begin{equation*}
L_\text{progress}(\psi)
\;=\; \mathbb{E}\!\left[ \bigl( \beta_\psi(s_t) - \tau(s_t) \bigr)^2 \right]
\;=\; \int_\mathcal{S} p(s')\, \mathbb{E}\!\left[ (\beta_\psi(s') - \tau(s_t))^2 \,\big|\, s_t = s' \right] \, \mathrm{d}s',
\end{equation*}
where $p(s')$ denotes the density of the state-visitation distribution. Differentiating with respect to $\beta_\psi(s)$, only the term at $s' = s$ contributes,
\begin{align*}
\frac{\partial L_\text{progress}}{\partial \beta_\psi(s)}
&\;=\; p(s)\, \frac{\partial}{\partial \beta_\psi(s)}\, \mathbb{E}\!\left[ (\beta_\psi(s) - \tau(s_t))^2 \,\big|\, s_t = s \right] \\
&\;=\; p(s)\, \mathbb{E}\!\left[ 2\, (\beta_\psi(s) - \tau(s_t)) \,\big|\, s_t = s \right] \\
&\;=\; 2\, p(s)\, \Bigl( \beta_\psi(s) - \mathbb{E}[\tau(s_t) \mid s_t = s] \Bigr).
\end{align*}
Identifying $\mathbb{P}(s_t = s) \equiv p(s)$ in the discrete or density-weighted sense, this gives the announced identity
\begin{equation*}
\frac{\partial L_\text{progress}}{\partial \beta_\psi(s)}
\;=\; 2\, \mathbb{P}(s_t = s)\, \Bigl( \beta_\psi(s) - \mathbb{E}[\tau(s_t) \mid s_t = s] \Bigr).
\end{equation*}
In particular, the partial derivative is positive when $\beta_\psi(s)$ overshoots the conditional mean target and negative when it undershoots, so a gradient descent step on $L_\text{progress}$ shrinks $\beta_\psi(s)$ in the first case and grows it in the second. Repeated application of this update therefore drives $\beta_\psi(s)$ monotonically toward $\mathbb{E}[\tau(s_t) \mid s_t = s]$ until equality is reached, in agreement with Proposition~\ref{prop:target_recovery}.
\end{proof}

\section{Proof of Lemma 3.1: Bounded Shaped Reward}
\label{app:proof_bounded}

\begin{proof}
By Equation~\eqref{eq:combined_reward},
\begin{equation*}
\bar r_t - R_t^E \;=\; \alpha\, \beta_\psi(s_t)\, I_t^{+,m}.
\end{equation*}
Taking absolute values and using $\alpha > 0$,
\begin{equation*}
\left| \bar r_t - R_t^E \right| \;=\; \alpha\, \left| \beta_\psi(s_t) \right|\, \left| I_t^{+,m} \right|.
\end{equation*}
By Definition~\ref{def:beta_network}, $\beta_\psi(s_t) \in [\beta_\text{min}, \beta_\text{max}]$ with $\beta_\text{min} > 0$, so $|\beta_\psi(s_t)| \le \beta_\text{max}$. By assumption, $|I_t^{+,m}| \le M$ almost surely. Combining these bounds,
\begin{equation*}
\left| \bar r_t - R_t^E \right| \;\le\; \alpha\, \beta_\text{max}\, M,
\end{equation*}
which establishes the claim.
\end{proof}

\section{Proof of Proposition 3.3: Computational Cost Comparison}
\label{app:proof_cost}

\begin{proof}
We measure cost in scalar floating-point operations and ignore lower-order terms.

\textbf{Meta-gradient training step.} A meta-gradient training step for a parameterized intrinsic reward function (LIRPG \cite{zheng2018lirpg}) consists of the following operations on a buffer of $T$ transitions:
(i) compute the policy gradient at the current parameters and apply $K$ inner update steps to obtain a perturbed policy $\theta'$, each of which costs $O(T\, |\theta|)$ to evaluate the policy gradient and $O(|\theta|)$ to apply the update; and
(ii) backpropagate the meta-objective evaluated at $\theta'$ through the chain of $K$ inner updates to obtain a gradient with respect to the intrinsic reward parameters. The backward pass through $K$ unrolled updates costs $O(K\, T\, |\theta|)$ in time and $O(K\, |\theta|)$ in memory because each inner step must be cached. Summing the forward and backward costs, the worst-case time complexity per buffer update is
\begin{equation*}
C_\text{meta} \;=\; O(K\, T\, |\theta|).
\end{equation*}

\textbf{CARE training step.} A CARE training step for the Beta Network on the same buffer of $T$ transitions consists of:
(i) one forward pass through the critic $V_\mu$ for each of the $T$ states to obtain TD errors $\delta_t$, at cost $O(T\, |\mu|)$, where $|\mu| \le |\theta|$;
(ii) construction of the target $\tau(s_t)$ from $|\delta_t|$, which is $O(T)$;
(iii) one forward pass through $\beta_\psi$ for each of the $T$ states, at cost $O(T\, |\psi|)$; and
(iv) one backward pass through $\beta_\psi$ to compute $\nabla_\psi L_\beta$, at cost $O(T\, |\psi|)$.
The dominant term is $O(T\, |\psi|)$ when $|\psi| \ge |\mu|$ in practice; otherwise the bound becomes $O(T\, |\mu|)$, which is still strictly smaller than $C_\text{meta}$ by a factor of $K$. We summarize the time complexity per buffer update as
\begin{equation*}
C_\text{CARE} \;=\; O(T\, |\psi|).
\end{equation*}

\textbf{Asymptotic ratio.} Taking the ratio,
\begin{equation*}
\frac{C_\text{meta}}{C_\text{CARE}} \;=\; \Theta\!\left( \frac{K\, |\theta|}{|\psi|} \right).
\end{equation*}
When $|\psi| \ll |\theta|$ and $K \ge 1$, this ratio is large, which is what the proposition asserts.
\end{proof}

\section{Proof of Corollary 3.1: Degeneration to Fixed Coefficient}
\label{app:proof_degeneration}

\begin{proof}
We treat the two regimes separately.

\textbf{(i) Cold-start regime.} Suppose $\max_t R_t^E \le 0$ within the rollout buffer. By Equation~\eqref{eq:beta_target}, the target reduces to the constant value $\tau(s_t) \equiv \beta_0$ for every $t$. Substituting into the progress loss in Equation~\eqref{eq:progress_loss},
\begin{equation*}
L_\text{progress}(\psi) \;=\; \mathbb{E}\!\left[ (\beta_\psi(s_t) - \beta_0)^2 \right].
\end{equation*}
This squared expectation is non-negative and equals zero if and only if $\beta_\psi(s) = \beta_0$ almost everywhere with respect to the state-visitation distribution. The regularizer $L_\text{reg}$ in Equation~\eqref{eq:reg_loss} is also non-negative and is minimized at the same point, since $L_\text{reg}(\psi) = \mathbb{E}[(\log \beta_\psi(s_t) - \log \beta_0)^2]$ vanishes precisely when $\beta_\psi(s) = \beta_0$. Therefore both terms of $L_\beta$ in Equation~\eqref{eq:beta_loss} are simultaneously minimized at $\beta_{\psi^\star}(s) \equiv \beta_0$.

\textbf{(ii) Strong regularization regime.} The progress loss $L_\text{progress}$ is non-negative and bounded above by $(\beta_\text{max} - \beta_\text{min})^2$ since both $\beta_\psi(s_t)$ and $\tau(s_t)$ take values in $[\beta_\text{min}, \beta_\text{max}]$. The regularizer $L_\text{reg}$ is non-negative. Let $\psi^\star_\lambda$ be any minimizer of $L_\beta$ at regularization weight $\lambda \equiv \lambda_\text{reg}$. By optimality, comparing against any reference $\psi_0$ that satisfies $\beta_{\psi_0}(s) \equiv \beta_0$ (for which $L_\text{reg}(\psi_0) = 0$),
\begin{equation*}
\gamma_p\, L_\text{progress}(\psi^\star_\lambda) \;+\; \lambda\, L_\text{reg}(\psi^\star_\lambda)
\;\le\; \gamma_p\, L_\text{progress}(\psi_0) \;+\; 0.
\end{equation*}
Rearranging,
\begin{equation*}
\lambda\, L_\text{reg}(\psi^\star_\lambda)
\;\le\; \gamma_p\, \bigl( L_\text{progress}(\psi_0) - L_\text{progress}(\psi^\star_\lambda) \bigr)
\;\le\; \gamma_p\, (\beta_\text{max} - \beta_\text{min})^2.
\end{equation*}
Dividing by $\lambda > 0$,
\begin{equation*}
L_\text{reg}(\psi^\star_\lambda)
\;\le\; \frac{\gamma_p\, (\beta_\text{max} - \beta_\text{min})^2}{\lambda}
\;\xrightarrow[\lambda \to \infty]{}\; 0.
\end{equation*}
Hence $\mathbb{E}[(\log \beta_{\psi^\star_\lambda}(s_t))^2] \to 0$, which implies $\log \beta_{\psi^\star_\lambda}(s_t) \to 0$ in $L^2$ under the state-visitation distribution. Equivalently, $\beta_{\psi^\star_\lambda}(s) \to \beta_0 = 1$ in $L^2$ on the support of that distribution.

\textbf{Reduction to fixed-coefficient form.} In either regime, substituting $\beta_\psi(s) = 1$ into Equation~\eqref{eq:combined_reward} yields
\begin{equation*}
\bar r_t \;=\; R_t^E \;+\; \alpha\, \cdot\, 1\, \cdot\, I_t^{+,m} \;=\; R_t^E \;+\; \alpha\, I_t^{+,m},
\end{equation*}
which is exactly the fixed-coefficient PPO+ICM update with intrinsic strength $\alpha$.
\end{proof}
